\chapter{Interdisciplinary Project Team Structure and Workflow}

This chapter describes the organisational structure of the Inscription and Manuscript Digitisation Project, detailing how the three disciplines—Computer Science (AI and software engineering), Electronics and Communication Engineering (signal processing and image quality analysis), and Industrial Engineering and Management (project management and impact assessment)—collaborate to deliver the Phase~1 image processing pipeline. Understanding the team structure, departmental roles, workflow methodologies, and technical infrastructure is essential for comprehending how specialised expertise from multiple domains converges toward the unified Phase~1 objective: transforming degraded historical scans into enhanced, binarised images ready for downstream processing in Phase~2.

\section{Interdisciplinary Team Composition and Roles}

The project is implemented by three discipline clusters, each with exclusive technical ownership of specific components and shared responsibility for system integration.

\subsection{Computer Science and Information Technology (CS/IT)}

The CS/IT team owns the end-to-end software pipeline from raw image ingestion to final binarised output.

\subsubsection{Responsibilities and Deliverables}

\begin{itemize}
	\item \textbf{Preprocessing stage:} Implementation of image normalisation, white balance correction, border cropping, and orientation handling (\texttt{src/preprocess.py}). This ensures consistent baseline quality before AI processing.

	\item \textbf{AI enhancement stage:} Development of super-resolution via Real-ESRGAN, non-local means denoising, DStretch colour decorrelation, and unsharp mask sharpening through artefact-specific processing chains (\texttt{src/enhance.py}). Owns all neural network inference and model weight management.

	\item \textbf{Binarisation stage:} Implementation of Sauvola adaptive thresholding, Otsu global thresholding, and OpenCV adaptive mean thresholding, plus post-processing morphological closing and noise blob removal (\texttt{src/binarise.py}). This stage is the Phase~1 endpoint; its binary PNG output is the deliverable passed to Phase~2 for character recognition.
\end{itemize}

\subsection{Electronics and Communication Engineering (ECE)}

The ECE team brings signal processing rigour to the image processing stages, characterising degradation mechanisms and designing custom filters for improved enhancement and binarisation quality.

\subsubsection{Contributions}

\textbf{Contribution 1 — Noise characterisation and modelling (\texttt{docs/noise\_analysis\_report.pdf}):}

ECE profiles noise types across all five artefact categories:
\begin{itemize}
	\item \textit{Gaussian noise}: Random pixel fluctuations from camera sensor thermal and photon shot noise.
	\item \textit{Salt-and-pepper noise}: Impulse noise from dust particles, scratches, and sensor read errors.
	\item \textit{JPEG compression artefacts}: Block boundaries and ringing from lossy JPEG encoding.
	\item \textit{Periodic noise}: Scanner line artefacts from mechanical imperfections.
\end{itemize}

Per-artefact noise profiles inform denoising strength selection: stone inscriptions (Gaussian + salt-and-pepper, denoising strength=10), palm leaf manuscripts (Gaussian + JPEG, strength=15), copper plates (specular reflectivity + Gaussian, strength=8).

\textbf{Contribution 2 — Custom filter design (\texttt{src/filters.py}):}

\begin{itemize}
	\item \textbf{Gabor filter bank:} Bank of 24 oriented spatial frequency filters (3 frequencies × 8 orientations) to decompose inscription images into oriented texture components, enabling character-background separation independent of intensity.

	\item \textbf{Directional edge enhancement:} Orientation-selective sharpening that enhances carving edges at specified angles (e.g., vertical strokes in palm leaf characters).

	\item \textbf{FFT-based periodic noise removal:} Frequency-domain notch filtering to suppress scanner artefacts while preserving text content (broadband).
\end{itemize}

\textbf{Contribution 3 — Colour channel and histogram analysis (\texttt{src/analysis.py}):}

Statistical characterisation of colour distributions:
\begin{itemize}
	\item Per-channel statistics: mean, standard deviation, skewness, kurtosis.
	\item DStretch colour space parameter tuning: empirical selection of LAB vs. YDS vs. YBK vs. LDS per artefact material.
	\item Before/after histogram visualisations for report figures.
\end{itemize}

\textbf{Contribution 4 — Image quality metrics (\texttt{src/metrics.py}):}

Development of the Phase~1 evaluation framework with four key metrics:

\begin{itemize}
	\item \textbf{Peak Signal-to-Noise Ratio (PSNR):} Quantifies fidelity of enhanced image vs. original. Target: $\geq 30$ dB.

	\item \textbf{Structural Similarity Index (SSIM):} Measures preservation of image structure (edges, contrast). Target: $\geq 0.85$.

	\item \textbf{Contrast-to-Noise Ratio (CNR):} Quantifies text-to-background contrast relative to background noise. Higher CNR indicates better text-background separability in the binarised output.

	\item \textbf{Edge sharpness:} Laplacian variance to verify that super-resolution maintains crisp character edges without over-smoothing.
\end{itemize}

Automatic per-image quality scoring is integrated into the pipeline and ensures every processed image is validated against objective criteria.

\subsection{Industrial Engineering and Management (IEM)}

The IEM team owns operational and organisational aspects: project scheduling, workflow optimisation, cost-benefit analysis, and stakeholder documentation.

\subsubsection{Contributions}

\textbf{Contribution 1 — Project management and scheduling (\texttt{docs/project\_plan.xlsx}):}

Weekly sprint planning, Gantt chart, task tracking, and risk register identifying:
\begin{itemize}
	\item \textit{Schedule risk}: Real-ESRGAN model weight acquisition and environment setup delays.
	\item \textit{Data risk}: Insufficient high-quality training datasets for Brahmi/Grantha scripts (Phase 2 risk).
	\item \textit{Technical risk}: GPU compute constraints for large-batch Real-ESRGAN processing.
	\item \textit{Quality risk}: Binarisation quality falling below threshold on heavily degraded artefacts.
\end{itemize}

\textbf{Contribution 2 — Digitisation workflow optimisation (\texttt{docs/workflow\_analysis.pdf}):}

Value stream mapping identifying the Phase~1 pipeline bottleneck:
\begin{itemize}
	\item Preprocessing (Stage 1): 0.5 seconds/image
	\item \textbf{AI Enhancement (Stage 2): 2.0 seconds/image (bottleneck — Real-ESRGAN GPU inference)}
	\item Binarisation (Stage 3): 0.3 seconds/image
\end{itemize}

Throughput analysis: current configuration processes approximately \textit{21 images/minute} on a single GPU. Parallelisation strategies (multi-GPU batch enhancement) are documented for scaling to institutional archive sizes.

\textbf{Contribution 3 — Cost-benefit and scalability analysis (\texttt{docs/cost\_benefit\_analysis.pdf}):}

Economic justification and scaling models:
\begin{itemize}
	\item \textbf{Cost per image}: GPU compute (\$0.09/image on AWS for enhancement), storage (\$0.03/image), total infrastructure \$0.12/image.
	\item \textbf{Manual enhancement alternative}: Expert manual image restoration costs \$3–8/image, requires 20–30 minutes per image.
	\item \textbf{ROI:} Breakeven at 15 images; automated enhancement becomes economically superior beyond that threshold.
	\item \textbf{Scalability}: 100 images (5 minutes), 10,000 images (8 hours), 1,000,000 images (33 days on 10 GPUs).
\end{itemize}

\textbf{Contribution 4 — Stakeholder documentation and impact assessment (\texttt{docs/user\_guide.pdf}, \texttt{docs/impact\_statement.pdf}):}

\begin{itemize}
	\item \textbf{User guide}: Non-technical instructions for museum curators, library staff, and heritage researchers to supply images and interpret enhanced and binarised outputs.

	\item \textbf{Impact statement}: Quantified benefits for institutional stakeholders (Archaeological Survey of India, Digital Library of India, eGangotri portal):
	\begin{itemize}
		\item Image legibility: 100,000+ inscriptions made visually accessible for study.
		\item Processing efficiency: 1,000 images enhanced and binarised per hour at full GPU capacity.
		\item Foundation for Phase~2: clean binary images eliminate the hardest preprocessing bottleneck for future character recognition work.
	\end{itemize}

	\item \textbf{Data governance policy}: Ownership, access control, CC BY 4.0 licensing, and compliance with heritage institution policies.
\end{itemize}

\section{Technical Infrastructure and Computing Environment}

\subsection{Hardware Configuration}

The development environment assumes:
\begin{itemize}
	\item \textbf{GPU:} NVIDIA GPU with CUDA 12.1 support (tested on RTX 4090, RTX A100). Real-ESRGAN inference: 2.0 seconds per image.
	\item \textbf{RAM:} 32 GB minimum (48 GB recommended). Batch processing of 10 images in memory.
	\item \textbf{Storage:} 500 GB SSD for code, models, and sample data. Cloud storage (AWS S3) for production archive.
\end{itemize}

\subsection{Software Dependencies and Versions}

\subsubsection{Image Processing}
\begin{itemize}
	\item OpenCV 4.9.0.80 — morphological operations, colour space conversions
	\item Pillow 10.3.0 — EXIF handling, image I/O
	\item scikit-image 0.23.2 — Sauvola binarisation, metric computation
	\item NumPy 1.26.4 — numerical array operations
\end{itemize}

\subsubsection{AI and Deep Learning}
\begin{itemize}
	\item PyTorch $\geq$ 2.0.0 — deep learning runtime
	\item BasicSR 0.3.0 / 1.4.2 — Real-ESRGAN models and inference
\end{itemize}

\subsection{Development Workflow and Version Control}

The project uses Git for version control with a branching strategy:
\begin{itemize}
	\item \texttt{main}: Stable, tested code (CI/CD pipeline runs on all merges)
	\item \texttt{develop}: Integration branch for feature merging
	\item \texttt{feature/*}: Per-feature branches (preprocessing, enhancement, binarisation, etc.)
\end{itemize}

Continuous integration via GitHub Actions validates all Phase~1 stages:
\begin{itemize}
	\item Python syntax and type checking (mypy)
	\item Unit tests (pytest) — 5 tests covering preprocessing and enhancement modules
	\item Image processing validation (PSNR threshold checks on known inputs)
\end{itemize}

\section{Operational Workflows and Methodologies}

\subsection{Development Lifecycle}

The project follows an iterative agile methodology with two-week sprints:

\begin{enumerate}
	\item \textbf{Sprint planning:} Team identifies bottlenecks from previous sprint (e.g., binarisation quality on heavily degraded stone samples), selects tasks (parameter tuning, new test cases).

	\item \textbf{Daily standup:} Asynchronous updates via Slack on blockers, progress, and interdisciplinary handoffs (e.g., ECE noise analysis inputs to CS denoising strength selection).

	\item \textbf{Code review:} All changes to \texttt{src/} require peer review before merging to develop. ECE contributions (filters, metrics) are reviewed for signal-processing correctness; IEM documentation is reviewed for completeness.

	\item \textbf{Integration testing:} Full Phase~1 pipeline tested on sample dataset (5 images per artefact type) to verify inter-stage consistency (e.g., binarised PNG is valid output with correct inversion polarity).

	\item \textbf{Sprint review:} Demo of enhanced and binarised outputs to stakeholders (simulated museum curators, heritage researchers) for feedback.

	\item \textbf{Sprint retrospective:} Team reflects on process improvements, scheduling accuracy, and technical debt.
\end{enumerate}

\subsection{Testing Strategy}

\begin{itemize}
	\item \textbf{Unit tests:} Per-function validation (e.g., \texttt{test\_preprocess.py} verifies CLAHE produces expected histogram shape).
	\item \textbf{Integration tests:} Multi-stage validation (e.g., preprocessing output is correctly ingested by enhancement and produces PSNR $\geq 30$\,dB).
	\item \textbf{Regression tests:} Known-good outputs are stored; new builds are verified against regression suite.
	\item \textbf{Visual inspection:} Team reviews enhanced and binarised outputs qualitatively; misaligned results trigger root-cause analysis.
\end{itemize}

\subsection{Interdisciplinary Collaboration Points}

\begin{itemize}
	\item \textbf{CS $\rightarrow$ ECE:} CS delivers enhanced image outputs; ECE computes quality metrics (PSNR, SSIM, CNR) on the test set and identifies underperforming artefact types.

	\item \textbf{ECE $\rightarrow$ CS:} ECE delivers custom filters (Gabor bank, FFT noise removal); CS integrates them as optional pipeline steps selectable per artefact type.

	\item \textbf{CS, ECE $\rightarrow$ IEM:} Technical teams provide weekly throughput metrics and quality scores; IEM updates the project plan and risk register.

	\item \textbf{IEM $\rightarrow$ CS, ECE:} IEM identifies bottlenecks (Real-ESRGAN is the slowest stage at 2.0\,s/image); technical teams propose solutions (batch GPU inference, model quantisation).
\end{itemize}

\section{Data Flow and Processing Architecture}

Table~\ref{tab:data_flow} summarises the Phase~1 data flow. Raw images from institutional archives (ASI, DLI, eGangotri) are ingested and routed sequentially through preprocessing, enhancement, and binarisation. Quality metrics are computed automatically after Stage~2, and all outputs are written to separate directories to preserve the non-destructive processing principle: original images are never overwritten.

\begin{table}[H]
\centering
\caption{Phase~1 data flow: stage inputs, processing, and outputs}
\label{tab:data_flow}
\begin{tabular}{|p{3.5cm}|p{5.5cm}|p{4cm}|}
\hline
\textbf{Input} & \textbf{Processing} & \textbf{Output} \\
\hline
Raw TIFF/JPG (degraded) & Stage~1: CLAHE, white balance, EXIF correction, border crop & Normalised JPG (quality=95) \\
\hline
Normalised JPG & Stage~2: Real-ESRGAN, denoising, DStretch, sharpening & Enhanced JPG \\
\hline
Enhanced JPG & Stage~3: Sauvola thresholding, morphological closing, blob removal & Binary PNG (\textbf{Phase~1 output}) \\
\hline
\end{tabular}
\end{table}

\section{Summary}

The interdisciplinary team structure uniquely positions the project to succeed by combining three complementary expertise domains. The CS/IT team leverages established deep learning frameworks and image processing libraries to deliver a robust, scalable three-stage pipeline. The ECE team ensures scientifically rigorous processing through noise characterisation, custom filter design, and objective quality validation. The IEM team manages project execution, optimises workflow efficiency, and communicates impact to institutional stakeholders. The technical infrastructure — GPU compute, Python image processing stack, and continuous integration — provides the computational foundation. Agile development methodology and iterative refinement ensure steady progress toward the Phase~1 goal of producing validated, high-quality binary images from degraded historical scans. Chapter~3 details the system architecture and design decisions that integrate these organisational capabilities into a cohesive technical solution.
